\begin{textsample}{POS Dim 6 – ac – Score 7.00 – ac/ac\_ef\_36.txt}  \label{ex:f6_pos_158}
5.
COMPETÊNCIAS GERAIS DA EDUCAÇÃO BÁSICA E DE ÁREA Com o objetivo de nortear escolas e professores na condução e aplicação dos componentes curriculares de forma interdisciplinar, as competências gerais da Base Nacional Comum Curricular estarão presentes em todos os níveis de ensino da educação básica, visando proporcionar aos alunos uma formação humana integral.
Essas competências se articulam na construção de conhecimentos, no desenvolvimento de habilidades e na formação de atitudes e valores, conforme preconizado na Lei de Diretrizes e Bases da Educação ( Lei nº 9.394/1996 ).
O trabalho com a língua inglesa envolve 6 competências específicas, as quais se acham articuladas com as 6 competências de linguagens, que por sua vez estão inseridas dentro das grandes competências da Base Nacional Comum Curricular, e devem garantir aos alunos o desenvolvimento de competências específicas.
COMPETÊNCIAS GERAIS DA BNCC PARA A EDUCAÇÃO BÁSICA 1.
Conhecimento - Valorizar e utilizar os conhecimentos historicamente construídos sobre o mundo físico, social, cultural e \textbf{digital} para entender e explicar a realidade, continuar aprendendo e colaborar para a construção de uma sociedade justa, democrática e inclusiva. 2.
Pensamento científico, crítico e criativo - Exercitar a curiosidade intelectual e recorrer à abordagem própria das ciências, incluindo a investigação, a reflexão, a análise crítica, a imaginação e a criatividade, para investigar causas, elaborar e \textbf{testar} hipóteses, formular e resolver \textbf{problemas} e criar soluções ( inclusive tecnológicas ) com base nos conhecimentos das diferentes áreas. 3.
Repertório cultural - Valorizar e fruir as diversas manifestações artísticas e culturais, das locais às mundiais, e participar de práticas diversificadas da produção \textbf{artístico-cultural}. 4.
Comunicação - Utilizar diferentes linguagens – verbal ( oral ou visual-motora, como Libras, e escrita ), corporal, visual, sonora e \textbf{digital} –, bem como conhecimentos das linguagens artística, matemática e científica, para se expressar e partilhar informações, experiências, ideias e sentimentos em diferentes contextos, além de produzir sentidos que levem ao entendimento mútuo. 5.
Cultura \textbf{digital} - Compreender, utilizar e criar \textbf{tecnologias} \textbf{digitais} de informação e comunicação de forma crítica, significativa, reflexiva e ética nas diversas práticas sociais ( incluindo as escolares ) para se comunicar, acessar e disseminar informações, produzir conhecimentos, resolver \textbf{problemas} e exercer protagonismo e autoria na vida pessoal e coletiva. 6.
Trabalho e projeto de vida - Valorizar a diversidade de saberes e vivências culturais, apropriar -se de conhecimentos e experiências que lhe possibilitem entender as relações próprias do mundo do trabalho e fazer escolhas alinhadas ao exercício da cidadania e ao seu projeto de vida, com liberdade, autonomia, consciência crítica e responsabilidade. 7.
Argumentação - Argumentar com base em fatos, dados e informações confiáveis, para formular, negociar e defender ideias, pontos de vista e decisões comuns que respeitem e promovam os direitos humanos, a consciência socioambiental e o consumo responsável em âmbito local, regional e global, com posicionamento ético em relação ao cuidado de si mesmo, dos outros e do planeta. 8.
Autoconhecimento e autocuidado - Conhecer -se, apreciar -se e cuidar de sua saúde física e emocional, compreendendo -se na diversidade humana e reconhecendo suas emoções e as dos outros, com autocrítica e capacidade para lidar com elas. 9.
Empatia e cooperação - Exercitar a empatia, o diálogo, a resolução de conflitos e a cooperação, fazendo -se respeitar e promovendo o respeito ao outro e aos direitos humanos, com acolhimento e valorização da diversidade de indivíduos e de grupos sociais, seus saberes, suas identidades, suas culturas e suas potencialidades, sem preconceitos de qualquer natureza. 10.
Responsabilidade e cidadania - Agir pessoal e coletivamente com autonomia, responsabilidade, flexibilidade, resiliência e determinação, tomando decisões com base em princípios éticos, democráticos, inclusivos, sustentáveis e solidários.
COMPETÊNCIAS DA BNCC DA ÁREA DE CONHECIMENTO 1.
Compreender as linguagens como construção humana, histórica, social e cultural, de natureza dinâmica, reconhecendo -as e valorizando -as como formas de significação da realidade e expressão de subjetividades e identidades sociais e culturais. 2.
Conhecer e explorar diversas práticas de linguagem ( artísticas, corporais e linguísticas ) em diferentes campos da atividade humana para continuar aprendendo, ampliar suas possibilidades de participação na vida social e colaborar para a construção de uma sociedade mais justa, democrática e inclusiva. 3.
Utilizar diferentes linguagens – verbal ( oral ou visual-motora, como Libras, e escrita ), corporal, visual, sonora e \textbf{digital} – para se expressar e partilhar informações, experiências, ideias e sentimentos em diferentes contextos e produzir sentidos que levem ao diálogo, à resolução de conflitos e à cooperação. 4.
Utilizar diferentes linguagens para defender pontos de vista que respeitem o outro e promovam os direitos humanos, a consciência socioambiental e o consumo responsável em âmbito local, regional e global, atuando criticamente frente a questões do mundo contemporâneo. 5.
Desenvolver o senso estético para reconhecer, fruir e respeitar as diversas manifestações artísticas e culturais, das locais às mundiais, inclusive aquelas pertencentes ao patrimônio cultural da humanidade, bem como participar de práticas diversificadas, individuais e coletivas, da produção \textbf{artístico-cultural}, com respeito à diversidade de saberes, identidades e culturas. 6.
Compreender e utilizar \textbf{tecnologias} \textbf{digitais} de informação e comunicação de forma crítica, significativa, reflexiva e ética nas diversas práticas sociais ( incluindo as escolares ), para se comunicar por meio das diferentes linguagens e mídias, produzir conhecimentos, resolver \textbf{problemas} e desenvolver projetos autorais e coletivos. ( BNCC, Brasil 2018 ). ( BRASIL, Base Nacional Comum Curricular, 2017 ).
COMPETÊNCIAS ESPECÍFICAS DO COMPONENTE INGLÊS - Identificar o lugar de si e o do outro em um mundo plurilíngue e multicultural, refletindo, criticamente, sobre como a aprendizagem da língua inglesa contribui para a inserção dos sujeitos no mundo globalizado, inclusive no que concerne ao mundo do trabalho. - Comunicar -se na língua inglesa, por meio do uso variado de linguagens e mídias impressas ou \textbf{digitais}, reconhecendo -a como ferramenta de acesso ao conhecimento, de ampliação das perspectivas e de possibilidades para a compreensão dos valores e interesses de outras culturas e para o exercício do protagonismo social. - Identificar similaridades e diferenças entre a língua inglesa e a língua materna/outras línguas, articulando -as a aspectos sociais, culturais e identitários, em uma relação intrínseca entre língua, cultura e identidade. - Elaborar repertórios linguístico-discursivos da língua inglesa, usados em diferentes países e por grupos sociais distintos dentro de um mesmo país, de modo a reconhecer a diversidade linguística como direito e valorizar os usos heterogêneos, híbridos e multimodais \textbf{emergentes} nas sociedades contemporâneas. - Utilizar novas \textbf{tecnologias}, com novas linguagens e modos de interação, para pesquisar, selecionar, compartilhar, posicionar -se e produzir sentidos em práticas de letramento na língua inglesa, de forma ética, crítica e responsável. - Conhecer diferentes patrimônios culturais, materiais e imateriais, difundidos na língua inglesa, com vistas ao exercício da fruição e da ampliação de perspectivas no contato com diferentes manifestações \textbf{artístico-culturais}.
Como forma de sintetizar e orientar, metodologicamente, o professor quanto ao currículo da disciplina de Língua Inglesa, apresentamos, a seguir, o Quadro Organizador Curricular com todos os objetivos de aprendizagem do 6º ao 9º ano, organizado por eixos \textbf{estruturantes} e conteúdos a serem trabalhados pelos professores, com o intuito de garantirmos direitos de aprendizagem de cada aluno, conforme o disposto na Base Nacional Comum Curricular.

% matched lemmas: artístico-cultural, digital, emergente, estruturante, problema, tecnologia, testar
\end{textsample}
